\documentclass[11pt]{article}
\usepackage[T1]{fontenc}
\usepackage{lmodern}
\usepackage[margin=1.1in]{geometry}
\usepackage{amsmath, amssymb, amsthm}
\usepackage{booktabs}
\usepackage{xcolor}
\definecolor{linknavy}{RGB}{25, 55, 125}
\usepackage[colorlinks=true, linkcolor=linknavy, citecolor=linknavy, urlcolor=linknavy]{hyperref}
\hypersetup{
  pdftitle={Cyclically 4-edge-connected cubic graphs of unbounded Petersen defect},
  pdfauthor={Felipe Santibáñez-Leal},
  pdfkeywords={Petersen coloring conjecture, normal edge-coloring, abnormal edges, sublinear approximation, cyclically 4-edge-connected cubic graphs, snarks, SAT, DRAT}
}
\usepackage{orcidlink}
\usepackage{fancyhdr}

\newtheorem{theorem}{Theorem}[section]
\newtheorem{proposition}[theorem]{Proposition}
\newtheorem{lemma}[theorem]{Lemma}
\newtheorem{corollary}[theorem]{Corollary}
\newtheorem{conjecture}[theorem]{Conjecture}
\theoremstyle{definition}
\newtheorem{definition}[theorem]{Definition}
\theoremstyle{remark}
\newtheorem{remark}[theorem]{Remark}

\newcommand{\PG}{P}
\newcommand{\pd}{\mathrm{pd}}
\newcommand{\ab}{\mathrm{ab}}
\newcommand{\Gfty}{G_{52}}
\newcommand{\Cut}{\mathrm{Cut}}

\newcommand{\doctype}{Preprint}
\newcommand{\docversion}{v0.02}
\newcommand{\docdate}{6 October 2026}
\newcommand{\versiondoi}{10.5281/zenodo.23196817}
\newcommand{\conceptdoi}{10.5281/zenodo.22847192}
\newcommand{\docrepo}{github.com/fsantibanezleal/CAOS\_RESEARCH}
\newcommand{\headerblock}{%
\begin{center}
\fbox{\parbox{0.94\textwidth}{\small\raggedright
\textbf{\doctype} \quad$\cdot$\quad \docversion, \docdate \quad$\cdot$\quad Licence CC BY 4.0\\[2pt]
CAOS Research open-problems programme \quad$\cdot$\quad problem \texttt{petersen-coloring}\\[2pt]
Version DOI \href{https://doi.org/\versiondoi}{\texttt{\versiondoi}}\\[2pt]
Concept DOI (always latest) \href{https://doi.org/\conceptdoi}{\texttt{\conceptdoi}}\\[2pt]
Not peer reviewed. Source code, data and computational records: \href{https://\docrepo}{\texttt{\docrepo}}
}}
\end{center}
\vspace{4pt}}

\title{Cyclically 4-Edge-Connected Cubic Graphs\\
of Unbounded Petersen Defect}
\author{Felipe Santib\'a\~nez-Leal\,\orcidlink{0000-0002-0150-3246}\\
\small CAOS Research program, \texttt{github.com/fsantibanezleal/CAOS\_RESEARCH}}
\date{\docdate}

\pagestyle{fancy}
\fancyhf{}
\fancyhead[L]{\small CAOS Research \doctype}
\fancyhead[R]{\small Unbounded Petersen Defect \textperiodcentered\ \docversion}
\fancyfoot[C]{\small \thepage}
\renewcommand{\headrulewidth}{0.4pt}
\fancypagestyle{plain}{\fancyhf{}\fancyfoot[C]{\small \thepage}\renewcommand{\headrulewidth}{0pt}}

\begin{document}
\maketitle
\headerblock

\begin{abstract}
The Petersen defect of a cubic graph is the least number of vertices at which an edge map into
the Petersen graph fails to send the vertex star onto a vertex star; it is zero exactly for the
graphs with a Petersen coloring, and it bounds from below the least number of abnormal edges of a
proper 5-edge-coloring. Mattiolo, Mazzuoccolo and Mkrtchyan conjectured that five statements are
equivalent: the Petersen coloring conjecture, and the existence of a sublinear bound on the number
of abnormal edges for all bridgeless, all 2-connected, all 3-connected, and all cyclically
4-edge-connected cubic graphs. After the refutation of the Petersen coloring conjecture in 2026,
the first two statements are false, ring and frame constructions show that the third and fourth
are false too, and the conjecture reduces to the last one. We show that it is false as well: for every $t \ge 2$ there is a cyclically 4-edge-connected cubic
graph of girth 5 on $102t$ vertices whose Petersen defect is at least $t$, so every proper
5-edge-coloring of it has at least $t$ abnormal edges. Hence all five statements are false and the
conjectured equivalence holds. The graphs are rings that alternate two 4-poles cut from a
52-vertex counterexample of Goedgebeur, Jooken, M\'a\v{c}ajov\'a, Mattiolo, Mazzuoccolo and
Ulyanov. The argument rests on a conservation law: modulo the cut space of the Petersen graph, a
4-pole whose vertices are all correctly mapped transmits the same class through both of its
connectors, and the class of a pair of labels is their distance in the line graph of the Petersen
graph. One 4-pole transmits only distance 1, the other only distance 0, so no two consecutive
blocks of a ring can both be free of bad vertices. The two transmitted sets and the non-colorability
of the smallest ring are certified by SAT refutations with DRAT proofs checked by
\texttt{drat-trim}, and the cyclic 4-edge-connectivity of all rings follows from that of the
four-block ring, verified by two independent exhaustive searches. The ring on 306 vertices has
Petersen defect exactly 3. For rings joined along 3-edge connectors, the class of a vertex modulo
the cut space is one of 64 cosets of the cut code of the Petersen graph, in six orbits under its
automorphism group, and the transfer relation of a 6-pole is block diagonal over these classes; so
the Petersen colorability of every ring built from a finite family of 6-poles, of every length, is
decided by a finite semigroup computation with a re-checkable certificate. Every ring of claws, of
Petersen superedges, or of both, is Petersen colorable. This covers the flower snarks, whose
junctions always receive three pairwise antipodal edges of the Petersen graph, and an infinite
family of cyclically 5-edge-connected snarks of girth 5. Splitting a 52-vertex counterexample along
a 6-edge cut gives two colorable 6-poles that conduct the same classes, so the classes alone do not
detect its non-colorability, while the transfer relation does.
\end{abstract}

\section{Introduction}

Let $\PG$ be the Petersen graph. A \emph{Petersen coloring} of a cubic graph $G$ is a map
$\sigma : E(G) \to E(\PG)$ such that for every vertex $v$ of $G$ the three edges at $v$ are sent
bijectively onto the three edges at some vertex of $\PG$. Jaeger conjectured that every bridgeless
cubic graph has one~\cite{Jaeger1988}, and showed that a cubic graph has a Petersen coloring if and
only if it has a \emph{normal} proper 5-edge-coloring, one in which every edge is poor (its two
end stars together see three colors) or rich (they see five)~\cite{Jaeger1985}. The conjecture was
refuted in August 2026 by Putman~\cite{Putman2026}, with a human-checkable proof by
Jooken~\cite{Jooken2026}; Goedgebeur, Jooken, M\'a\v{c}ajov\'a, Mattiolo, Mazzuoccolo and
Ulyanov~\cite{GJMMMU2026} found two counterexamples on 52 vertices and cyclically
4-edge-connected counterexamples of every even order at least 60.

An edge of a proper 5-edge-coloring that is neither poor nor rich is \emph{abnormal}; write
$\ab(G)$ for the least number of abnormal edges over all proper 5-edge-colorings of $G$. Mattiolo,
Mazzuoccolo and Mkrtchyan~\cite{MMM2021} proved that no proper 5-edge-coloring has exactly one
abnormal edge and made the following conjecture.

\begin{conjecture}[\cite{MMM2021}, Conjecture~3]\label{conj:mmm}
The following statements are equivalent: (a) the Petersen coloring conjecture holds; there is a
sublinear function $f$ with $\ab(G) \le f(|V(G)|)$ for (b) all bridgeless, (c) all 2-connected,
(d) all 3-connected, respectively (e) all cyclically 4-edge-connected cubic graphs $G$.
\end{conjecture}

They proved that (a) and (b) are equivalent, and that on each of the last three classes a
sublinear bound exists if and only if $\ab$ is bounded by 5, 7 and 9
respectively~\cite[Theorems~1--4]{MMM2021}. After the disproof, (a) and (b) are false. Rings of
counterexamples joined through 2-edge cuts, and cubic frames whose vertices are replaced by
counterexamples minus a vertex, have one bad vertex in every block, so (c) and (d) are false as
well~\cite{Audit2026}. The same work shows that, given a proposition of~\cite{MMM2021} on
cyclic joins, (e) is equivalent to every cyclically 4-edge-connected cubic graph admitting a map
into $E(\PG)$ with at most two bad vertices, and it records that in fifteen cyclically
4-edge-connected counterexamples every pair of adjacent vertices is critical. The conjecture
therefore turned on statement (e). We settle it; the proof does not use that proposition.

\begin{theorem}\label{thm:main}
For every $t \ge 2$ there is a cyclically 4-edge-connected cubic graph $R_t$ of girth~5 on $102t$
vertices such that every map $E(R_t) \to E(\PG)$ has at least $t$ vertices at which the star
condition fails. Consequently every proper 5-edge-coloring of $R_t$ has at least $t$ abnormal
edges.
\end{theorem}

\begin{corollary}\label{cor:e}
No sublinear function bounds the number of abnormal edges on cyclically 4-edge-connected cubic
graphs. Hence all five statements of Conjecture~\ref{conj:mmm} are false, and the conjecture
holds.
\end{corollary}

The graphs $R_t$ are rings of $2t$ blocks that alternate two 4-poles, $A$ and $B$, both cut from
the first 52-vertex counterexample $\Gfty$ of~\cite{GJMMMU2026}: $A$ is $\Gfty$ with two
independent edges removed, and $B$ is $\Gfty$ with two adjacent vertices removed. The proof uses a
conservation law. Modulo the cut space of $\PG$, a 4-pole whose vertices all satisfy the star
condition transmits the same class through its two connectors (Lemma~\ref{lem:conduction}), and
the class carried by a pair of edges is the distance of their labels in the line graph of $\PG$
(Lemma~\ref{lem:distance}). The 4-pole $A$ transmits only distance~1 and $B$ only distance~0
(Proposition~\ref{prop:sets}), so two consecutive blocks of a ring cannot both be free of bad
vertices (Theorem~\ref{thm:ring}). Two small restoration arguments (Lemma~\ref{lem:restore}) show
that each of the two computed restrictions is enough on its own. The rings are cyclically
4-edge-connected because the four-block ring is (Lemma~\ref{lem:c4ec}).

The rings have 4-edge cuts, and every known counterexample to the Petersen coloring conjecture
does; whether a cyclically 5-edge-connected one exists is Problem~11 of~\cite{GJMMMU2026}.
Section~\ref{sec:sixpoles} examines the natural analogue of the construction for that class, rings
of 6-poles joined along 3-edge connectors. The class of a connector modulo the cut space is still
conserved, but on the 6-poles of the Petersen graph, the flower snarks $J_5$, $J_7$ and the
dodecahedron it never separates two blocks. The exact transfer relation is block diagonal over the
64 classes, and the colorability of every ring of a finite family of blocks reduces to a finite
semigroup. For the claw and the Petersen superedge every ring is colorable
(Theorem~\ref{thm:ringsY}).

\section{Preliminaries}\label{sec:prelim}

Graphs are finite. For a vertex $v$, $\partial(v)$ is the set of edges at $v$; for a vertex set
$S$, $\delta(S)$ is the set of edges with exactly one end in $S$. A cubic graph is
\emph{cyclically 4-edge-connected} if every edge cut whose two sides both contain a cycle has at
least four edges.

\paragraph{Maps into the Petersen graph.} For a cubic graph $G$ and a map $\sigma : E(G) \to
E(\PG)$, a vertex $v$ is \emph{good} if $\sigma$ maps $\partial(v)$ bijectively onto the star of a
vertex of $\PG$, and \emph{bad} otherwise. The \emph{Petersen defect} $\pd(G)$ is the least number of
bad vertices over all maps; $\pd(G) = 0$ exactly when $G$ has a Petersen coloring. For $v$ let
$\chi_v \in \mathbb{F}_2^{E(\PG)}$ be the sum of the indicator vectors of the three labels at $v$,
and let $\Cut(\PG)$ be the cut space of $\PG$. The star of a vertex is a cut, so $\chi_v \in
\Cut(\PG)$ when $v$ is good; for every map the sum of $\chi_v$ over all vertices is zero, because
every edge is counted at both ends. The parity theorem of~\cite{Audit2026} follows: the number of
bad vertices is never exactly one, so every graph without a Petersen coloring has $\pd \ge 2$.

\paragraph{Abnormal edges.} We use the Kneser model of $\PG$: its vertices are the 2-subsets of
$\{1, \dots, 5\}$, two of them adjacent when disjoint, and the edge $\{X, Y\}$ has the color that
lies outside $X \cup Y$, so the three edges at $X$ carry the three colors outside $X$.

\begin{lemma}[\cite{Audit2026}]\label{lem:pdab}
$\pd(G) \le \ab(G)$ for every cubic graph $G$.
\end{lemma}

\begin{proof}
Given a proper 5-edge-coloring $c$, let $X_v$ be the pair of colors missing at $v$, and let the
\emph{view} of an edge $e = uv$ from $u$ be the edge of $\PG$ at $X_u$ of color $c(e)$. If $e$ is
poor then $X_u = X_v$; if it is rich then $X_u \cap X_v = \emptyset$ and both avoid $c(e)$; in both
cases the two views coincide. Map every normal edge to its common view and every abnormal edge to
its view from one chosen end. A vertex all of whose edges carry its own views is good, so the bad
vertices are among the ends not chosen of the abnormal edges.
\end{proof}

\paragraph{Multipoles.} A \emph{4-pole} is obtained from a cubic graph by removing some vertices
and edges so that exactly four \emph{dangling ends} remain: an edge removed with one end kept, or
kept at one end only, becomes a dangling end at the kept vertex. The dangling ends are grouped into
two \emph{connectors} of two ends each. A map of a 4-pole assigns a label to every edge and to
every dangling end; the star condition is required at the vertices of the 4-pole.

\section{Classes of pairs of labels}\label{sec:classes}

The line graph of $\PG$ has diameter 3: from a fixed edge, four edges are at distance 1, eight at
distance 2 and two at distance 3. We write $d(x, y)$ for this distance.

\begin{lemma}\label{lem:smallcuts}
The cuts of $\PG$ with at most four edges are the empty cut, the ten stars, and the fifteen sets of
four edges adjacent to one edge.
\end{lemma}

\begin{proof}
$\PG$ is 3-edge-connected, cyclically 5-edge-connected and of girth~5. If one side $S$ of a cut
induces a forest with $k$ components, the cut has $3|S| - 2(|S| - k) = |S| + 2k$ edges, which is 3
only for a single vertex and 4 only for two adjacent vertices; if both sides contain a cycle, the
cut has at least five edges.
\end{proof}

\begin{lemma}\label{lem:distance}
Let $x, y, x', y'$ be edges of $\PG$. If $\mathbf{1}_x + \mathbf{1}_y$ and $\mathbf{1}_{x'} +
\mathbf{1}_{y'}$ lie in the same class modulo $\Cut(\PG)$, then $d(x, y) = d(x', y')$.
\end{lemma}

\begin{proof}
The class of $\mathbf{1}_x + \mathbf{1}_y$ is zero exactly when $x = y$, since a vector of weight 2
is not a cut. If $x \ne y$, $x' \ne y'$ and the two sums have the same class, the symmetric
difference of $\{x, y\}$ and $\{x', y'\}$ is a cut of even weight at most 4. By
Lemma~\ref{lem:smallcuts} it is empty, so that the pairs coincide, or it consists of the four
edges adjacent to an edge $st$ of $\PG$. Then either each pair consists of two edges at the same
end of $st$, and both pairs are at distance 1, or each pair has one edge at $s$ and one at $t$,
and both pairs are at distance 2 because $\PG$ has no triangle.
\end{proof}

\begin{lemma}[conduction]\label{lem:conduction}
Let $M$ be a 4-pole with connectors $\{p, q\}$ and $\{r, s\}$, and let $\sigma$ be a map of $M$ in
which every vertex of $M$ is good. Then $d(\sigma(p), \sigma(q)) = d(\sigma(r), \sigma(s))$.
\end{lemma}

\begin{proof}
The sum of $\chi_v$ over the vertices of $M$ is a sum of stars, hence a cut. An edge of $M$ is
counted at both of its ends, so the sum equals $\mathbf{1}_{\sigma(p)} + \mathbf{1}_{\sigma(q)} +
\mathbf{1}_{\sigma(r)} + \mathbf{1}_{\sigma(s)}$. The two connectors therefore carry the same class,
and Lemma~\ref{lem:distance} gives the claim.
\end{proof}

For a 4-pole $M$ with its connectors, the \emph{transmitted set} $D(M) \subseteq \{0, 1, 2, 3\}$ is
the set of distances $d(\sigma(p), \sigma(q))$ over all maps $\sigma$ of $M$ with every vertex
good. By Lemma~\ref{lem:conduction} it does not depend on which connector is read. Because the
automorphism group of $\PG$ is transitive on ordered pairs of edges at a given distance, $d \in
D(M)$ can be decided by one formula in which the labels of $p$ and $q$ are fixed to one such pair.

\section{Rings of 4-poles}\label{sec:rings}

Let $A$ and $B$ be 4-poles, each with a first and a second connector. The ring $R_t(A, B)$ consists
of $t$ copies of $A$ and $t$ copies of $B$ placed alternately around a cycle, the second connector
of each block joined by two edges to the first connector of the next block. It is cubic. The two
edges that join consecutive blocks form a connector of both blocks.

\begin{theorem}\label{thm:ring}
If $D(A) \cap D(B) = \emptyset$, then every map of $R_t(A, B)$ has bad vertices in at least $t$ of
its $2t$ blocks. In particular $\ab(R_t(A, B)) \ge \pd(R_t(A, B)) \ge t$.
\end{theorem}

\begin{proof}
If two consecutive blocks had only good vertices, Lemma~\ref{lem:conduction} applied to each of
them would place the distance of their common connector in $D(A) \cap D(B)$. So the blocks without
bad vertices are pairwise non-consecutive, and at most $t$ of the $2t$ blocks are among them. The
last inequality is Lemma~\ref{lem:pdab}.
\end{proof}

\begin{lemma}[restoration]\label{lem:restore}
Let $G$ be a cubic graph without a Petersen coloring.
\begin{enumerate}
\item[\textup{(a)}] If $ab$ and $cd$ are independent edges and $M = G - ab - cd$ has connectors
$\{a, b\}$, $\{c, d\}$ (the dangling ends at these vertices), then $0 \notin D(M)$.
\item[\textup{(b)}] If $uv$ is an edge, $u_1, u_2$ are the other neighbors of $u$ and $v_1, v_2$ those of
$v$, and $M = G - \{u, v\}$ has connectors $\{u_1, u_2\}$, $\{v_1, v_2\}$, then $1 \notin D(M)$.
\end{enumerate}
\end{lemma}

\begin{proof}
(a) Distance 0 means equal labels on $a$ and $b$, and then by Lemma~\ref{lem:conduction} equal labels
on $c$ and $d$; putting these labels on $ab$ and $cd$ gives a Petersen coloring of $G$. (b) If
the labels at $u_1, u_2$ are at distance 1, they share a vertex of $\PG$ whose third edge $z$
satisfies $\mathbf{1}_{\sigma(uu_1)} + \mathbf{1}_{\sigma(uu_2)} \equiv \mathbf{1}_z$ modulo
$\Cut(\PG)$; likewise the labels at $v_1, v_2$ give an edge $z'$. By Lemma~\ref{lem:conduction} the
two classes agree, so $\mathbf{1}_z + \mathbf{1}_{z'}$ is a cut and $z = z'$. Labelling $uv$ by $z$
makes $u$ and $v$ good, which gives a Petersen coloring of $G$.
\end{proof}

So for $A = G - ab - cd$ and $B = G' - \{u, v\}$ with $G$, $G'$ without Petersen colorings, the sets
$D(A)$ and $D(B)$ are disjoint as soon as $D(A) \subseteq \{1\}$ or $D(B) \subseteq \{0\}$.

\begin{lemma}\label{lem:c4ec}
Let $G$ be a cyclically 4-edge-connected cubic graph, $A = G - ab - cd$ and $B = G - \{u, v\}$ as in
Lemma~\ref{lem:restore}, with $A$ joined through its connector $\{c, d\}$ to the connector $\{u_1,
u_2\}$ of the next $B$, and $B$ through $\{v_1, v_2\}$ to the connector $\{a, b\}$ of the next $A$.
If $R_2(A, B)$ has no cycle-separating edge cut with at most three edges, then every $R_t(A, B)$
with $t \ge 2$ is cyclically 4-edge-connected.
\end{lemma}

\begin{proof}
$G$ is 3-edge-connected, so $A$ is connected; $G$ is 3-connected, so $B$ is connected. Add to a
block a vertex $s^*$ joined to the two ends of its first connector and a vertex $t^*$ joined to
the two ends of its second connector. For $B$ this gives $G - uv$ with $u, v$ renamed $s^*, t^*$,
which has two edge-disjoint $s^*$-$t^*$ paths because $G$ is 3-edge-connected. For $A$, an edge
separating $s^*$ from $t^*$ would either separate $\{a, b\}$ from $\{c, d\}$ in $G$ by itself,
contradicting 3-edge-connectivity, or be an edge $s^*a$ (or $s^*b$, $t^*c$, $t^*d$) whose removal
leaves $b$ (respectively $a$, $d$, $c$) without a path to the other connector, contradicting the
connectivity of $A$. Consequently a ring carries four edge-disjoint paths between any two of its
blocks, two along each arc.

Let $K$ be a cycle-separating cut of $R_t$ with at most three edges and sides $S$, $S'$. By the
four paths, no block lies entirely in $S$ while another lies entirely in $S'$; say that every block
not split by $K$ lies in $S$. Every split block contains an edge of $K$, being connected, so at
most three blocks are split and $S'$ lies within them. A cycle in $S'$ lies in a maximal run $Y$ of
consecutive split blocks; $Y$ has at most three blocks and is bounded on both sides by blocks in
$S$. The edges leaving $S' \cap V(Y)$ belong to $K$ and form a cycle-separating cut whose size
depends only on $Y$ and $S' \cap V(Y)$. The ring $R_2 = ABAB$ contains every alternating run of at
most three blocks with at least one further block around it, so it contains the same cut,
contradicting the hypothesis.
\end{proof}

\section{The construction}\label{sec:construction}

Let $\Gfty$ be the 52-vertex counterexample of~\cite{GJMMMU2026} (House of Graphs entry 57244), with
its vertices numbered $0, \dots, 51$ as in the edge list of the repository (SHA-256 of the sorted
edge list \texttt{27db5d3b\ldots}); it is cyclically 4-edge-connected with girth~5 and has no
Petersen coloring~\cite{GJMMMU2026, Audit2026}. Let
\[
A = \Gfty - \{0,3\} - \{1,9\}, \qquad B = \Gfty - \{2, 7\},
\]
where $\{0, 3\}$, $\{1, 9\}$ are independent edges and $\{2, 7\}$ is an edge, with connectors as in
Lemma~\ref{lem:restore}: $\{0, 3\}$ and $\{1, 9\}$ for $A$, the two other neighbors of $2$ and the
two other neighbors of $7$ for $B$. Let $R_t = R_t(A, B)$, joined as in Lemma~\ref{lem:c4ec}; it has
$52t + 50t = 102t$ vertices.

\begin{proposition}[machine-verified]\label{prop:sets}
$D(A) = \{1\}$ and $D(B) = \{0\}$.
\end{proposition}

\begin{proof}
For each distance $d$ one propositional formula states that the 4-pole has a map with every vertex
good and the first connector labelled by a fixed pair at distance $d$. The satisfiable formulas
give maps that an independent checker accepts; the others are refuted with DRAT proofs accepted by
\texttt{drat-trim} (Table~\ref{tab:sets}). Distance 0 for $A$ and distance 1 for $B$ are also
excluded by Lemma~\ref{lem:restore}.
\end{proof}

\begin{table}[h]\centering\small
\begin{tabular}{lcccc}
\toprule
4-pole & $d = 0$ & $d = 1$ & $d = 2$ & $d = 3$ \\
\midrule
$A = \Gfty - \{0,3\} - \{1,9\}$ & refuted (7.7 s) & map (0.5 s) & refuted (2.7 s) & refuted (2.1 s) \\
$B = \Gfty - \{2,7\}$ & map (0.2 s) & refuted (2.8 s) & refuted (2.2 s) & refuted (2.1 s) \\
\bottomrule
\end{tabular}
\caption{Transmitted distances of the two 4-poles. ``Map'': a map with every vertex good and the
first connector at distance $d$, accepted by the checker (the second connector is then at the same
distance). ``Refuted'': unsatisfiable, with a DRAT proof accepted by \texttt{drat-trim}. Times are
solver times.}\label{tab:sets}
\end{table}

\begin{proposition}[machine-verified]\label{prop:rings}
The graphs $R_1$ and $R_2$ are simple cubic graphs of girth~5 without cycle-separating edge cuts of
at most three edges, and the four edges joining one block to the rest form a cycle-separating
cut, so their cyclic edge connectivity is exactly~4. $R_1$ has no Petersen coloring, and
$\pd(R_2) = 2$.
\end{proposition}

\begin{proof}
The cut searches are exhaustive: for every pair of edges, the bridges of the graph without that
pair are computed, and every edge set of size at most three found in this way is tested for
leaving two sides that contain cycles; the same routine applied to $\Gfty$, to a ring of two copies
of $\Gfty$ opened at one edge, and to the cubic graph $K_4$ with its vertices replaced by copies of
$\Gfty$ minus a vertex returns no cut, 2-edge cuts, and 3-edge cuts respectively, in agreement with
an enumeration of all edge sets of size at most three. That enumeration, run on $R_1$ and $R_2$
themselves, finds no cycle-separating cut either. The Petersen coloring formula of
$R_1$ is refuted with a DRAT proof accepted by \texttt{drat-trim}. For $R_2$, a map with exactly two
bad vertices, both in the two copies of $B$, is accepted by the checker; two is the least possible
by the parity theorem.
\end{proof}

\begin{proof}[Proof of Theorem~\ref{thm:main}]
By Proposition~\ref{prop:sets}, $D(A) \cap D(B) = \emptyset$, so Theorem~\ref{thm:ring} gives
$\ab(R_t) \ge \pd(R_t) \ge t$. By Proposition~\ref{prop:rings} and Lemma~\ref{lem:c4ec}, $R_t$ is
cyclically 4-edge-connected for every $t \ge 2$. A cycle of length at most four would lie within
two consecutive blocks and their joining edges, so it would occur in $R_2$, which has girth 5.
\end{proof}

\begin{proposition}[machine-verified]\label{prop:r3}
$\pd(R_3) = 3$.
\end{proposition}

\begin{proof}
The lower bound is Theorem~\ref{thm:ring}; it was also tested directly on the 306 vertices of $R_3$:
the formula stating that some map has at most two bad vertices is unsatisfiable, with a DRAT proof
accepted by \texttt{drat-trim} ($5{,}802$ seconds of solving and $2{,}265$ seconds of checking).
For the upper bound, the star condition is relaxed at one vertex in each of the three copies of
$B$, the vertex at the position (within its copy) of one of the two bad vertices of a two-defect
map of $R_2$; the relaxed formula is satisfiable (169 seconds), and the checker finds exactly these
three vertices bad.
\end{proof}

\begin{proof}[Proof of Corollary~\ref{cor:e}]
The graphs $R_t$ have $n = 102t$ vertices and $\ab(R_t) \ge n/102$, so no sublinear function bounds
$\ab$ on cyclically 4-edge-connected cubic graphs, and (e) is false; for $t \ge 10$ they also
violate the bound 9 of~\cite[Theorem~4]{MMM2021}. Statements (a) to (d) are false, as recalled in
the introduction, so the five statements are equivalent.
\end{proof}

\section{6-poles: classes and transfer relations}\label{sec:sixpoles}

\paragraph{The 64 classes.} Write $Q = \mathbb{F}_2^{E(\PG)}/\Cut(\PG)$, and call the class in $Q$ of
the sum of the indicator vectors of a multiset of labels the class of that multiset. The cut space
is a binary $[15, 9, 3]$ code (its nonzero words of weight at most 4 are those of
Lemma~\ref{lem:smallcuts}), so $Q$ has 64 elements. A vertex is good exactly when the class of its
three labels is $0$ and they are distinct. The automorphism group of $\PG$ acts on $Q$ with six
orbits (Table~\ref{tab:classes}); we call them $\mathsf{0}$, $\mathsf{E}$, $\mathsf{D}_3$,
$\mathsf{D}_2$, $\mathsf{T}_1$, $\mathsf{T}_2$. The orbit $\mathsf{T}_1$ is the class of the
all-ones vector: removing from $E(\PG)$ the twelve edges at the four 2-subsets containing a fixed
element of $\{1, \dots, 5\}$ leaves three pairwise antipodal edges (at distance 3 in the line graph),
and these five \emph{antipodal triples} partition $E(\PG)$.

\begin{table}[h]\centering\small
\begin{tabular}{lccl}
\toprule
orbit & classes & least weight & elements of least weight \\
\midrule
$\mathsf{0}$ & 1 & 0 & the cut space, containing the ten stars \\
$\mathsf{E}$ & 15 & 1 & one edge; also two edges at distance 1 \\
$\mathsf{D}_3$ & 15 & 2 & two edges at distance 3 \\
$\mathsf{D}_2$ & 30 & 2 & two edges at distance 2 (two such pairs per class) \\
$\mathsf{T}_1$ & 1 & 3 & the five antipodal triples \\
$\mathsf{T}_2$ & 2 & 3 & triples of pairwise distance-2 edges (ten per class) \\
\bottomrule
\end{tabular}
\caption{The orbits of $\mathrm{Aut}(\PG)$ on the 64 classes modulo the cut space (computed and
checked against Lemma~\ref{lem:distance}).}\label{tab:classes}
\end{table}

\paragraph{6-poles.} A \emph{6-pole} has six dangling ends, split into two connectors of three. The
proof of Lemma~\ref{lem:conduction} shows that a map with every vertex good puts the same class on
both connectors, and the \emph{conducted set} $D(M) \subseteq Q$ is the set of these classes; it
is invariant under $\mathrm{Aut}(\PG)$, hence a union of orbits. Theorem~\ref{thm:ring} holds
verbatim for rings of 6-poles joined along their connectors by three edges, and such rings can be
cyclically 5-edge-connected, since two arcs of a ring are separated by six edges. For a graph $G$
without a Petersen coloring and non-adjacent vertices $u$, $w$, the 6-pole $G - u - w$ with the
ends at $u$ and at $w$ as connectors does not conduct $\mathsf{0}$, by the argument of
Lemma~\ref{lem:restore}: class $\mathsf{0}$ at the ends of $u$ means three labels forming a star.

\begin{proposition}[machine-verified]\label{prop:core}
Let $H$ be the Petersen graph, the flower snark $J_5$ or $J_7$, or the dodecahedron. Every 6-pole
$H - u - w$ ($u$, $w$ non-adjacent, connectors at $u$ and at $w$), and, for $H \ne J_7$, every
6-pole $H - e_1 - e_2 - e_3$ ($e_1, e_2, e_3$ pairwise disjoint, any split of the six ends) conducts
$\mathsf{0}$, $\mathsf{E}$ and $\mathsf{D}_2$; the set $\{\mathsf{0}, \mathsf{E}, \mathsf{D}_2\}$
itself is conducted by $\PG - u - w$.
\end{proposition}

\begin{proof}
One formula for each 6-pole up to automorphisms of $H$, each split and each orbit, with the class of
the first connector fixed to a representative of the orbit (1,890 splits, 11,340 formulas). The
satisfiable ones give maps accepted by the checker, the others are refuted with DRAT proofs
accepted by \texttt{drat-trim}; none of the four graphs has a cycle-separating cut of at most four
edges.
\end{proof}

So no two of these 6-poles have disjoint conducted sets, and Theorem~\ref{thm:ring} cannot show
that a ring of them has no Petersen coloring. The exact obstruction is the following relation.

\paragraph{Transfer relations.} A \emph{block} is a 6-pole whose connectors $L$ and $R$ are ordered.
Its transfer relation is the set $T_M$ of pairs $(\sigma(L), \sigma(R)) \in E(\PG)^3 \times E(\PG)^3$
over all maps $\sigma$ with every vertex good. In a ring of blocks $M_1, \dots, M_t$, end $j$ of $R_i$
is joined to end $\pi_i(j)$ of $L_{i+1}$ (indices modulo $t$). Every vertex of the ring lies in one
block, so a Petersen coloring of the ring is a labelling of the $3t$ joining edges accepted by
every block: the ring is Petersen colorable if and only if the boolean matrix product $T_{M_1}
\Pi_1 T_{M_2} \Pi_2 \cdots T_{M_t} \Pi_t$ has a nonzero diagonal entry, where $\Pi_i$ is the
permutation matrix of $\pi_i$.

\begin{proposition}\label{prop:sectors}
For $q \in Q$ let $S_q$ be the set of ordered label triples of class $q$.
\begin{enumerate}
\item[\textup{(a)}] $T_M \subseteq \bigcup_{q \in Q} S_q \times S_q$, and every $\Pi_i$ maps each $S_q$
onto itself, so every product of transfer relations and junctions is block diagonal over $Q$.
\item[\textup{(b)}] The diagonal blocks of such a product over two classes of one orbit are conjugate,
so they have the same trace.
\end{enumerate}
Hence a ring is Petersen colorable if and only if, for some orbit, the product of the restrictions
to $S_q \times S_q$, for one fixed class $q$ of that orbit, has a nonzero diagonal. These six
\emph{sector matrices} have orders $60, 67, 48, 48, 30, 60$.
\end{proposition}

\begin{proof}
(a) is the conduction property, and the class of a triple does not depend on the order of its
entries. (b) An automorphism $\alpha$ of $\PG$ maps a map $\sigma$ of $M$ to the map $\alpha \circ
\sigma$, so $(x, y) \in T_M$ if and only if $(\alpha x, \alpha y) \in T_M$; $\alpha$ acts on labels
and $\Pi_i$ on positions, so they commute, and $\alpha$ conjugates the block over $q$ to the block
over $\alpha(q)$. The sector orders add up to $60 + 15 \cdot 67 + 15 \cdot 48 + 30 \cdot 48 + 30 + 2
\cdot 60 = 15^3$.
\end{proof}

\begin{corollary}\label{cor:semigroup}
Let $\mathcal{F}$ be a finite set of blocks and $\mathcal{G}$ the set of six-tuples of sector
matrices of $T_M \Pi$, for $M \in \mathcal{F}$ in either orientation and $\pi$ any bijection. Every
ring of at least two blocks of $\mathcal{F}$ is Petersen colorable if and only if every product of at
least two elements of $\mathcal{G}$ has a nonzero diagonal in some sector. This holds if there is a
set $\mathcal{C}$ of six-tuples of boolean matrices such that every product of two elements of
$\mathcal{G}$ contains an element of $\mathcal{C}$ entrywise, every product $cg$ with $c \in
\mathcal{C}$ and $g \in \mathcal{G}$ contains an element of $\mathcal{C}$, and every element of
$\mathcal{C}$ has a nonzero diagonal in some sector; and such a set exists whenever the condition
holds (take all products).
\end{corollary}

\begin{proof}
The first statement is Proposition~\ref{prop:sectors}. For the second, boolean products are
monotone under entrywise containment, so by induction on the length every product of at least two
elements of $\mathcal{G}$ contains an element of $\mathcal{C}$, and its diagonal is nonzero where
that element's is. The set of all products is finite, which gives the converse.
\end{proof}

A set $\mathcal{C}$ as in Corollary~\ref{cor:semigroup} is a certificate that anyone can check by
products of boolean matrices of order at most 67. Entries of computed sector matrices are maps
accepted by the checker, so a computed matrix can only miss entries; missing entries can only
lower diagonals, and a certificate built from computed matrices is a certificate for the exact
ones.

\paragraph{Claws and Petersen superedges.} The \emph{claw} $Y$ is the block with a center joined to
three vertices, each of which carries one dangling end in $L$ and one in $R$. The \emph{Petersen
superedge} $S$ is $\PG - u - w$ with $u$, $w$ at distance 2, $L$ the three ends at $u$ and $R$ the
three ends at $w$; it is the superedge $(P_{10})_{u,v}$ of Sedlar and \v{S}krekovski~\cite{SS2023}.
The rings of claws contain the flower snarks of Isaacs~\cite{Isaacs1975}: $J_k$ for odd $k$ is the
ring of $k$ claws with the identity at $k - 1$ junctions and a transposition at the last one.

\begin{theorem}[machine-verified]\label{thm:ringsY}
Every ring of at least two blocks, each a claw or a Petersen superedge in either orientation, joined
through arbitrary bijections, has a Petersen coloring. A ring of Petersen superedges always has one
whose joining edges carry classes of the orbit $\mathsf{E}$.
\end{theorem}

\begin{proof}
The sector matrices of $Y$ and $S$ were computed by enumerating the models of a SAT formula
projected to the dangling ends, each model accepted by the checker, and recomputed by exhaustive
enumeration of maps and by a second encoding with a second solver, with identical results. The set
of all products of at least two generators was closed exhaustively for $\{Y\}$, $\{S\}$ and
$\{Y, S\}$: it has $30$, $19{,}005$ and $116{,}463$ elements, all with a nonzero diagonal; for
$\{S\}$ every element has a nonzero diagonal in the sector $\mathsf{E}$. Certificates in the sense of
Corollary~\ref{cor:semigroup}, kept minimal under containment, have $12$, $276$ and $2{,}454$
elements and were re-verified by a separate program with a different matrix product.
\end{proof}

\begin{corollary}[flower snarks]\label{cor:flower}
For every odd $k \ge 3$, every Petersen coloring of $J_k$ maps the three edges joining two
consecutive claws onto an antipodal triple of $\PG$.
\end{corollary}

\begin{proof}
The products $(Y\,\mathrm{id})^{k-1}(Y\,\tau)$ are periodic in $k$ from $k = 5$ with period 2,
and for odd $k$ only the sector $\mathsf{T}_1$ has a nonzero diagonal (for even $k$ only
$\mathsf{E}$). By conduction all junctions of a coloring carry one class, which is therefore
$\mathsf{T}_1$, and the 30 ordered triples of class $\mathsf{T}_1$ are the orderings of the five
antipodal triples.
\end{proof}

Hägglund and Steffen showed that flower snarks have Petersen colorings~\cite{HS2014};
Corollary~\ref{cor:flower} describes all of them at the junctions.

\begin{proposition}\label{prop:superedgerings}
Let $S^t$ be the ring of $t$ Petersen superedges with the identity at every junction.
\begin{enumerate}
\item[\textup{(a)}] For $t \ge 3$, $S^t$ is a simple cubic graph, and it is 3-edge-colorable if and
only if $t$ is even.
\item[\textup{(b)}] For $t \ge 5$, $S^t$ has girth 5 and is cyclically 5-edge-connected.
\end{enumerate}
So for every odd $t \ge 5$, $S^t$ is a cyclically 5-edge-connected snark of girth 5 on $8t$ vertices,
and by Theorem~\ref{thm:ringsY} it has a Petersen coloring.
\end{proposition}

\begin{proof}
(a) The same reduction applies to 3-edge-colorings, with the 27 color triples of a connector as
states: the boundary relation of $S$, enumerated exhaustively, has powers of period 2, and the odd
powers have a zero diagonal. (b) An exhaustive search finds no cycle-separating cut with at most
four edges in $S^5$, whose girth is 5. $\PG - u - w$ is connected, and restoring $u$ and $w$ gives
$\PG$, which is 3-edge-connected, so a ring carries six edge-disjoint paths between any two blocks.
The argument of Lemma~\ref{lem:c4ec}, with cuts of at most four edges, at most four split blocks
and runs of at most four consecutive blocks (all of which occur in $S^5$), shows that $S^t$ has no
cycle-separating cut with at most four edges for $t \ge 5$, and a cycle of length at most four would
lie in such a run.
\end{proof}

\begin{remark}[the classes are not enough]\label{rem:control}
$\Gfty$ has a 6-edge cut with sides of 18 and 34 vertices. For each of the ten splits of the cut
into two triples, the two sides are 6-poles with Petersen colorings whose conducted sets share
$\mathsf{0}$, $\mathsf{E}$ and $\mathsf{D}_2$, so no argument by classes separates them; their ring
of length 2 is $\Gfty$, and the product of their sector matrices has a zero diagonal in all six
sectors, as it must. In a further test, 203 rings (the flower snarks $J_5$, $J_7$, $J_9$ and 200
random words of length 2 to 6 over 81 blocks, with 18 to 132 vertices) were decided both by the
sector matrices and by a SAT solver on the graph, with the same answer every time.
\end{remark}

\section{Remarks and questions}\label{sec:remarks}

\emph{Other 4-poles.} The computation behind Proposition~\ref{prop:sets} was carried out for a
wider set of 4-poles (computation EXP-012). The Petersen graph minus two adjacent vertices, and
the Petersen graph minus two edges at distance 2 or 3, transmit distance 1 under every pairing of
their dangling ends. For $\Gfty$ minus two adjacent vertices, with the pairing of
Lemma~\ref{lem:restore}(b), the transmitted set is $\{0\}$ for eleven of the fourteen edge orbits
of the automorphism group of $\Gfty$, $\{0, 3\}$ for two and $\{0, 2\}$ for one; it never contains
1, as Lemma~\ref{lem:restore}(b) requires. Under the two other pairings every such 4-pole
transmits distance 1. Among the 4-poles $\Gfty - ab - cd$ with the connectors of
Lemma~\ref{lem:restore}(a), 282 of the 546 edge pairs examined have $D = \{1\}$ (computation
EXP-010). Many rings with the property of Theorem~\ref{thm:ring} can therefore be built from
$\Gfty$ alone.

\emph{Criticality.} A pair of vertices of a graph without a Petersen coloring is \emph{critical}
if some map has exactly these two vertices bad. In every one of fifteen cyclically
4-edge-connected counterexamples examined in~\cite{Audit2026} every pair of adjacent vertices is
critical. By Theorem~\ref{thm:ring} the two bad vertices of a map of $R_2$ lie in two
non-consecutive blocks, so two adjacent vertices of one block of $R_2$ form a non-critical pair,
and in $R_t$ with $t \ge 3$ no pair is critical. Universal 2-criticality is therefore not a
property of the cyclically 4-edge-connected class.

\emph{Questions.} The rings have 4-edge cuts, and no cyclically 5-edge-connected counterexample to
the Petersen coloring conjecture is known~\cite[Problem~11]{GJMMMU2026}. A ring of 6-poles without
a Petersen coloring needs blocks whose semigroup in Corollary~\ref{cor:semigroup} contains a
product with a zero diagonal in every sector; which 6-poles without small internal cuts have this
property? Is the Petersen defect bounded on cyclically 5-edge-connected cubic graphs? What is the
largest constant $c$ such that cubic graphs with $\pd \ge c\,n$ exist in the cyclically
4-edge-connected class? The rings give $c \ge 1/102$.

\section*{Data availability}

{\raggedright
The edge list of $\Gfty$, the encoders, the checkers, the ring construction, the cut searches and
all results, with the SHA-256 digests of every formula and proof, are in the repository
\url{https://github.com/fsantibanezleal/CAOS_RESEARCH}, under
\path{problems/combinatorics/petersen-coloring/}. The computations of this paper are archived in the
folder \path{experiments/EXP-012-conduction-rings} of that directory (runners \path{run.py} and
\path{run_ring.py}; library modules \path{code/pcclib/poles.py} and \path{code/pcclib/rings.py}),
with the distance sets of~\cite{Audit2026} in \path{experiments/EXP-010-sublinear-approximation}.
The upper bound of Proposition~\ref{prop:r3} is in the same folder (\path{run_r3_upper.py}).
Section~\ref{sec:sixpoles} is archived in \path{experiments/EXP-013-six-pole-charges} (conducted
sets) and \path{experiments/EXP-014-transfer-semigroups} (sector matrices, closures, certificates
and their checker \path{check_certificate.py}, the control on $\Gfty$), with the library modules
\path{code/pcclib/charges.py} and \path{code/pcclib/transfer.py} and the probes
\path{code/probes/superedge_ring_*.py}; the theory notes are in
\path{context/2026-10-06-charges.md}. Proofs were produced by CaDiCaL~\cite{CaDiCaL} and checked with
\texttt{drat-trim}~\cite{DRATtrim}.
\par}

\begin{thebibliography}{99}
\bibitem{Jaeger1988} F.~Jaeger, Nowhere-zero flow problems, in: L.~W. Beineke, R.~J. Wilson
(eds.), Selected Topics in Graph Theory 3, Academic Press, London, 1988, 71--95.
\bibitem{Jaeger1985} F.~Jaeger, On five-edge-colorings of cubic graphs and nowhere-zero flow
problems, Ars Combinatoria 20-B (1985), 229--244.
\bibitem{Putman2026} B.~Putman, A 112-vertex counterexample to the Petersen Coloring Conjecture,
Zenodo v1.1.0 (2026), \href{https://doi.org/10.5281/zenodo.21845291}{doi:10.5281/zenodo.21845291};
arXiv:2608.10012.
\bibitem{Jooken2026} J.~Jooken, A human-checkable proof of the 112-vertex counterexample to the
Petersen coloring conjecture, arXiv:2608.10028v2 (2026).
\bibitem{GJMMMU2026} J.~Goedgebeur, J.~Jooken, E.~M\'a\v{c}ajov\'a, D.~Mattiolo, G.~Mazzuoccolo,
S.~Ulyanov, Disproving the Petersen Coloring Conjecture: theoretical analysis and an infinite
family of counterexamples, arXiv:2608.10028v4 (2026).
\bibitem{HS2014} J.~H\"agglund, E.~Steffen, Petersen-colorings and some families of snarks, Ars
Math. Contemp. 7 (2014), 161--173.
\bibitem{Isaacs1975} R.~Isaacs, Infinite families of nontrivial trivalent graphs which are not Tait
colorable, Amer. Math. Monthly 82 (1975), 221--239.
\bibitem{SS2023} J.~Sedlar, R.~\v{S}krekovski, Normal 5-edge-coloring of some snarks
superpositioned by the Petersen graph, arXiv:2305.05981 (2023).
\bibitem{MMM2021} D.~Mattiolo, G.~Mazzuoccolo, V.~Mkrtchyan, On sublinear approximations for the
Petersen coloring conjecture, Bull. Inst. Combin. Appl. 92 (2021), 78--90; arXiv:2104.09241.
\bibitem{Audit2026} F.~Santib\'a\~nez-Leal, Berge--Fulkerson covers, cycle double covers, flows and
exact normality defects of the first counterexamples to the Petersen coloring conjecture,
preprint v0.06 (2026), \href{https://doi.org/10.5281/zenodo.23195171}{doi:10.5281/zenodo.23195171}.
\bibitem{CaDiCaL} A.~Biere, T.~Faller, K.~Fazekas, M.~Fleury, N.~Froleyks, F.~Pollitt, CaDiCaL 2.0,
in: Computer Aided Verification (CAV 2024), Lecture Notes in Comput. Sci. 14681, Springer, 2024,
133--152.
\bibitem{DRATtrim} N.~Wetzler, M.~J.~H. Heule, W.~A. Hunt Jr., DRAT-trim: efficient checking and
trimming using expressive clausal proofs, in: Theory and Applications of Satisfiability Testing
(SAT 2014), Lecture Notes in Comput. Sci. 8561, Springer, 2014, 422--429.
\end{thebibliography}

\end{document}
